\section{Unified Finite Converse}\label{sec:converse}
%==============================================================================


This section isolates the one-shot obstruction that the later graph-capacity theory will lift and iterate. The cleanest zero-error formulation starts from the joint observation-tag map itself. The deterministic finite foundation below isolates one obstruction that later reappears as colorability, strong powers, and asymptotic capacity. The first theorem gives the basic criterion; the remaining statements record the same obstruction in combinatorial, entropic, and finite-error coordinates.

\begin{theorem}[Pair-injectivity characterization]\label{thm:pair-injective}
Let $X$ range over a finite latent alphabet, let $Y$ be the deterministic observation transcript, and let $T$ be a finite auxiliary tag. Exact zero-error recovery from $(Y,T)$ is possible if and only if the pair map $x \mapsto (Y(x),T(x))$ is injective on the ambiguity class under consideration.
\end{theorem}
\leanmeta{\LHrng{OBS}{1}{2}}

\begin{proof}
Let $\phi(x):=(Y(x),T(x))$.

\emph{Necessity.} If $\phi$ is not injective, then there exist distinct latent states $x_1\neq x_2$ with $\phi(x_1)=\phi(x_2)$. Any deterministic decoder $D$ therefore receives the same input on both states and must satisfy $D(\phi(x_1))=D(\phi(x_2))$. Since $x_1\neq x_2$, the common output cannot equal both latent states, so $D$ errs on at least one of them.

\emph{Sufficiency.} If $\phi$ is injective, define a decoder on the image of $\phi$ by $D(y,t)=\phi^{-1}(y,t)$, with arbitrary extension off the image. Injectivity makes $\phi^{-1}$ single-valued on the image, and for every latent state $x$ we then have $D(Y(x),T(x))=x$. Thus exact zero-error recovery is possible.
\end{proof}

\begin{proposition}[Confusability formulation]\label{thm:confusability-converse}
Fix a deterministic observation transcript and an $L$-bit side-information tag. If $K$ latent states induce the same observation transcript, then exact zero-error decoding on that ambiguity class requires $K$ distinct tag outcomes. Equivalently, a $K$-way ambiguity class forms a confusability clique whose size cannot exceed the available tag alphabet.
\end{proposition}
\leanmeta{\LH{OBS5}}

\begin{proof}
Apply Theorem~\ref{thm:pair-injective} on an ambiguity class on which the observation is constant. Then injectivity of $(Y,T)$ reduces to injectivity of the tag coordinate alone on that class. Since an $L$-bit tag provides at most $2^L$ outcomes, the clique size cannot exceed $2^L$.
\end{proof}

\paragraph{Global observation-tag budget.}
If the observation alphabet has size $O$ and the auxiliary tag alphabet has size $T$, then exact zero-error recovery on a latent alphabet of size $K$ requires
\[
K \le O\,T.
\]
\leanmeta{\LH{OBS4}}

\begin{proof}[Proof sketch]
Theorem~\ref{thm:pair-injective} requires injectivity of the map into the product alphabet $\mathcal Y \times \mathcal T$. A product alphabet of size $O\,T$ cannot support an injective image of more than $O\,T$ latent states.
\end{proof}

\paragraph{Fiber-level injectivity and cardinality.}
Fix an observation value $y$. On the observation fiber
\[
\mathcal F_y=\{x : Y(x)=y\},
\]
exact zero-error recovery forces the tag map to be injective. Consequently,
\[
|\mathcal F_y| \le |\mathcal T|.
\]
\leanmeta{\LH{OBS3}}

\begin{proof}[Proof sketch]
Restrict Theorem~\ref{thm:pair-injective} to $\mathcal F_y$. Since $Y$ is constant on the fiber, injectivity of the pair map reduces to injectivity of the tag coordinate alone. An injective map from $\mathcal F_y$ into the tag alphabet immediately gives $|\mathcal F_y| \le |\mathcal T|$.
\end{proof}

\begin{proposition}[Finite counting converse]\label{thm:fano-converse}
Let $F$ range over $K$ possible latent values. Suppose exact recovery is attempted from a fixed observation transcript together with an $L$-bit side-information tag. Then
\[
K \le 2^L.
\]
Equivalently, exact zero-error recovery of $K$ ambiguous states requires at least $\log_2 K$ bits of side information.
\end{proposition}
\leanmeta{\LH{CIA3}}

\begin{proof}
This is Theorem~\ref{thm:confusability-converse} specialized to a single $K$-way ambiguity class.
\end{proof}

The counting bound is the coarsest form of the obstruction. The next step is to weight the same argument by source mass. Exact recovery still enforces injectivity on successful confusability classes, but now the question is how much entropy can remain once the source is partitioned into success and failure branches.

\paragraph{Conditional-entropy formulation.}
Let $X$ be uniform on a $K$-way ambiguity class and let $Y$ be the deterministic observation transcript. If $Y$ is constant on that class, then $H(X\mid Y)=\log_2 K$. In particular, any tag-observation pair $(Y,T)$ that resolves the class exactly must satisfy
\[
H(X\mid Y,T)=0
\qquad\text{and}\qquad
H(T\mid Y)\ge \log_2 K .
\]
\leanmeta{\LH{PRB47}, \LH{PRB55}}

\begin{proof}[Proof sketch]
If $Y$ is constant on the class, conditioning on $Y$ leaves the uniform prior unchanged, hence $H(X\mid Y)=\log_2 K$. Exact recovery from $(Y,T)$ forces $H(X\mid Y,T)=0$, so the side information carried by $T$ must close the entire $\log_2 K$ gap.
\end{proof}

\paragraph{Mass-weighted clique entropy bound.}
Let $S$ be a success set for a zero-error decoder, and suppose $S$ forms a confusability clique under the observation transcript. If $p_S$ denotes the total probability mass of $S$, then the entropy carried by the successful states obeys
\[
\sum_{x\in S} p(x)\log_2 \frac{1}{p(x)}
\;\le\;
h_2(p_S) + p_S \log_2 |\mathcal T|,
\]
where $h_2$ is binary entropy and $|\mathcal T|$ is the tag alphabet size.
\leanmeta{\LH{PRB55}}

\begin{proof}[Proof sketch]
Inside a confusability clique, exact recovery on the success set still forces tag injectivity. The successful states therefore contribute at most $\log_2 |\mathcal T|$ bits per unit success mass. Separating success from failure adds the Bernoulli term $h_2(p_S)$, which yields the stated bound.
\end{proof}

In particular, on any exact success set inside a confusability clique, the restricted observation-tag map is injective by Theorem~\ref{thm:pair-injective}, so the same entropy inequality applies directly to the successful branch rather than only to the whole clique.\leanmeta{\LH{OBS1}, \LH{PRB55}}

\paragraph{Decoder-output and gap formulations.}
In the same deterministic finite model, the obstruction can also be expressed as output-entropy and finite-budget gap constraints: any deterministic decoder output $\hat X$ satisfies
\[
H(\hat X)\le H(Y,T)\le H(X),
\]
and the observation-tag entropy and decoded-output entropy are bounded by their corresponding finite alphabet ceilings, so the deficits
\[
\log_2 |\mathcal Y\times \mathcal T|-H(Y,T),
\qquad
\log_2 |\widehat{\mathcal X}|-H(\hat X)
\]
are nonnegative and vanish only in the uniform saturation cases formalized in the supplement.
\leanmeta{\LHrng{PRB}{73}{74}, \LHrng{PRB}{82}{83}, \LH{PRB90}, \LHrng{PRB}{93}{95}}

\begin{proof}[Proof sketch]
The decoder output is a deterministic function of $(Y,T)$, so the deterministic data-processing statement below gives $H(\hat X)\le H(Y,T)$. The pair entropy is itself bounded by the source entropy and by the logarithm of the available pair alphabet, which yields the nonnegative gap constraints.
\end{proof}

\paragraph{Deterministic data processing.}
Let $\kappa$ be any deterministic coarsening of the observation-tag pair $(Y,T)$. Then
\[
H(\kappa(Y,T)) \le H(Y,T).
\]
In particular, since any deterministic decoder output $\hat X$ is a coarsening of $(Y,T)$,
\[
H(\hat X)\le H(Y,T)\le \log_2 |\mathcal Y\times \mathcal T|.
\]
\leanmeta{\LH{PRB68}, \LH{PRB81}, \LH{PRB83}}

\begin{proof}[Proof sketch]
Deterministic coarsening merges atoms of the distribution on $(Y,T)$ and therefore cannot increase entropy. A decoder is one such coarsening. The ceiling follows because $(Y,T)$ takes values in a finite alphabet of size $|\mathcal Y\times\mathcal T|$.
\end{proof}

This theorem is the structural reason the output and gap formulations are not independent embellishments of the counting converse. Once the observation is coarsened, every derived representation inherits the same obstruction through deterministic data processing rather than escaping it.\leanmeta{\LH{PRB68}, \LH{PRB81}, \LH{PRB83}, \LH{PRB90}}

\begin{proposition}[Equivalence viewpoint]\label{thm:equivalence-viewpoint}
The confusability, counting, conditional-entropy, decoder-output, and finite-gap statements above are different normal forms of the same deterministic finite zero-error obstruction: a surviving $K$-way ambiguity class requires a budget of at least $\log_2 K$ bits to be resolved exactly.
\end{proposition}
\leanmeta{\LH{OBS1}, \LH{OBS5}, \LH{CIA3}, \LH{PRB47}, \LH{PRB55}, \LH{PRB68}, \LH{PRB81}, \LH{PRB83}, \LH{PRB90}}

\begin{proof}
Pair injectivity is the structural core. The confusability and counting theorems express the injectivity requirement combinatorially; the conditional-entropy and weighted-entropy theorems express it in source-mass coordinates; the decoder-output and gap theorems express it through deterministic coarsening and finite alphabet ceilings; and the finite-error theorem is the relaxed version obtained by allowing a nonzero failure branch. Each theorem therefore measures the same failure of exact isolation in a different coordinate system.
\end{proof}

\subsection{Finite-error extension}\label{sec:finite-error-extension}

The main body of the paper studies exact zero error. The same deterministic finite model also supports a budgeted finite-error extension, which should be read as a deterministic finite analogue of the classical Fano line of argument rather than as a separate asymptotic coding theorem~\cite{fano1961transmission,witsenhausen1975conditional,cover2006elements}.

\begin{proposition}[Budgeted finite-error converse]\label{thm:finite-error-budgeted}
Let $P_e$ denote the decoder error probability on a finite source of size $K$, observed through a deterministic transcript alphabet of size $O$ together with a tag alphabet of size $T$. Then the decoded-output entropy obeys
\[
H(\hat X)
\;\le\;
h_2(P_e) + (1-P_e)\log_2 (OT) + P_e \log_2 (K-1).
\]
\end{proposition}
\leanmeta{\LH{PRB85}, \LH{PRB87}}

\begin{proof}
Partition the source into success and failure events. On the success branch, the decoded output is constrained by the observation-tag budget and contributes at most $\log_2(OT)$. On the failure branch, the decoder can still output at most one of the remaining $K-1$ alternatives. The Bernoulli split between the two branches contributes the binary-entropy term. Equivalently,
\[
H(\hat X)=H(\hat X \mid \text{success/failure}) + H(\text{success/failure})
\]
is bounded by combining the support bound on each branch with the binary entropy of the branch variable itself.
\end{proof}

The zero-error theorems are the $P_e=0$ boundary of this inequality. In that regime the success branch occupies all probability mass, the Bernoulli term vanishes, the failure contribution disappears, and the budget collapses back to the unified finite converse above.
